\documentclass[12pt, a4paper]{ctexart}

% --- 1. 数学核心宏包 📐 ---
\usepackage{amsmath, amssymb, amsthm, amsfonts}
% --- 2. 页面与超链接设置 🔗 ---
\usepackage{geometry}  % 保持默认宽边距，不额外设置则默认为标准边距
\usepackage[colorlinks=true, linkcolor=black]{hyperref}
% --- 3. 其它必要包 ---
\usepackage{enumitem}
% --- 4. 封面信息预设 ---
\title{Note For INTRODUCTION TO LINEAR ALGEBRA}
\author{Jaehaerys Wang}
\date{\today}

% --- 正文区 ---
\begin{document}

% 自动生成封面
\maketitle
\newpage

% 自动生成目录
\tableofcontents
\newpage

\section{Chapter 1: Introduction to Vectors}

% 写法示例
\newpage
\section*{Examples}
\subsection*{逻辑结构与列表}
\begin{itemize} % label = $\bullet$ 可以自定义无序列表符号
    \item 这是一个无序列表项。
    \item 这是另一个无序列表项。
    \begin{itemize}
        \item 这是一个嵌套的无序列表项。
        \item 这是另一个嵌套的无序列表项。
    \end{itemize}
\end{itemize}
\begin{enumerate} % label = \arabic* 可以自定义有序列表符号
    \item 这是一个有序列表项。
    \begin{itemize}[nosep] % 嵌套无序列表，紧凑格式
        \item 这是一个嵌套的无序列表项。
        \item 这是另一个嵌套的无序列表项。
    \end{itemize}
    \item 这是另一个有序列表项。
    \begin{enumerate}[label=性质 \arabic*:] % 嵌套有序列表
        \item 这是一个嵌套的有序列表项。
        \item 这是另一个嵌套的有序列表项。
    \end{enumerate}
\end{enumerate}
\subsection*{数学语境与基础记号}
% 在 LaTeX 中，排版数学内容主要分为两种模式：行内模式 (Inline Math)：公式嵌入在文字中间，使用 $ ... $。例如：设 $n$ 维向量 $\mathbf{x}$。行间模式 (Display Math)：公式独立成行且居中，使用 \[ ... \]。
$\mathbb{R} \in \forall \subseteq \mathbb{R}^n \space \mathbf{v} \vec{v} $ \\
\subsection*{矩阵与方程组排版}
% 所有的矩阵环境都遵循相同的内部逻辑：
% 使用 & 来分隔同一行中的不同列。
% 使用 \\ 来换行，分隔不同的行。
% 无括号 (no brackets)
$ \begin{matrix} 1&2\\3&4 \end{matrix} $
% 圆括号 (parentheses)
$ \begin{pmatrix} 1&2\\3&4 \end{pmatrix} $
% 方括号 (brackets)
$ \begin{bmatrix} 1&2\\3&4 \end{bmatrix} $
% 单竖线 (vertical bars，常用于行列式)
$ \begin{vmatrix} 1&2\\3&4 \end{vmatrix} $
% 双竖线 (double vertical bars，常用于范数)
$ \begin{Vmatrix} 1&2\\3&4 \end{Vmatrix} $ \\
% 省略号的使用
% 水平省略号 (horizontal ellipsis)
$ \begin{bmatrix} 1&2&\cdots&n\\0
    &0&\cdots&0\\\vdots&\vdots&\ddots&\vdots\\0&0&\cdots&0 \end{bmatrix} $
% 垂直省略号 (vertical ellipsis)
$ \begin{bmatrix} 1&2&\cdots&n\\0
    &0&\cdots&0\\\vdots&\vdots&\ddots&\vdots\\0&0&\cdots&0 \end{bmatrix} $
% 对角线省略号 (diagonal ellipsis)
$ \begin{bmatrix} 1&2&\cdots&n\\0
    &0&\cdots&0\\\vdots&\vdots&\ddots&\vdots\\0&0&\cdots&0 \end{bmatrix} $
% 方程组与大括号
$ \begin{cases} x+y=1\\ x-y=2 \end{cases} $
% 高斯消元与步骤对齐
\begin{align*}
    \begin{bmatrix}
        1 & 2 \\
        2 & 5
    \end{bmatrix} 
    &\xrightarrow{R_2 - 2R_1} % 在箭头处对齐
    \begin{bmatrix}
        1 & 2 \\
        0 & 1
    \end{bmatrix} \\
    &\xrightarrow{R_1 - 2R_2} % 下一行的箭头也会跟上面对齐
    \begin{bmatrix}
        1 & 0 \\
        0 & 1
    \end{bmatrix}
\end{align*}

\end{document}